% ==============================================================================
% T0/FFGFT-Theorie: Dokument 241
% Titel: Die zwei Schichten – Verhältnisse und Umrechnung.
%        Eine pädagogische Einführung in die methodische Architektur der FFGFT
% Autor: Johann Pascher
% Datum: Mai 2026
% Bezug: Dokumente 014 (nat./SI), 015 (NatEinheitenSystematik), 044 (Feinstruktur-
%        konvention), 052 (Massen-Elimination), 133 (Fraktale Korrektur),
%        134 (c als Konvention), 011 (Feinstruktur-Herleitung), 205 (Narrativ),
%        159 (Konzertsaal/Modenstruktur), 161 (Ising), 167 (LiNbO3), 173 (Quantenpunkt),
%        185 (Einbettungspreis), 230 (Hilbertraum-Brücke)
% ==============================================================================

\section*{Vorbemerkung}

Dieses Dokument richtet sich an Leserinnen und Leser, die in die methodische
Architektur der FFGFT einsteigen möchten, ohne sich zuvor durch den vollen
mathematischen Apparat zu arbeiten. Es wiederholt keine Ableitungen; es
erklärt eine Lese-Reihenfolge.

Der Inhalt ist nicht neu. Die zentrale Aussage --- dass FFGFT in zwei klar
getrennten Schichten arbeitet --- steht bereits an verschiedenen Stellen
der Dokumentenreihe (vor allem in Dok.~014, 015, 044 und 134). Sie steht
dort aber als technische Klarstellung neben anderen technischen Aussagen,
nicht als architektonisches Prinzip mit eigener Stimme. Genau dies soll
hier nachgeholt werden.

Wer die Theorie zum ersten Mal liest, läuft sonst Gefahr, an einer Stelle
hängen zu bleiben, die sich auflöst, sobald man die zwei Schichten erkennt.
Wer sie nicht erkennt, sieht in jedem fraktalen Korrekturfaktor einen freien
Parameter und in jeder Naturkonstante eine eigenständige Größe --- und
übersieht damit den eigentlichen Punkt: dass die FFGFT mit genau einem
freien Parameter $\xi_0 = 4/30000$ auskommt, und dass alle übrigen Größen
entweder reine Verhältnisse sind oder über eine geometrisch bestimmte
Skalenbrücke entstehen.

\section{Der Ausgangspunkt: das Konstantenproblem}

Die übliche Schulphysik präsentiert eine Liste von ``fundamentalen
Naturkonstanten'':
\begin{itemize}\setlength\itemsep{0pt}
  \item die Lichtgeschwindigkeit $c$,
  \item die reduzierte Planck-Konstante $\hbar$,
  \item die Gravitationskonstante $G$,
  \item die Boltzmann-Konstante $k_B$,
  \item die elektrische Feldkonstante $\varepsilon_0$,
  \item die Feinstrukturkonstante $\alpha \approx 1/137$,
  \item die Elementarladung $e$,
\end{itemize}
und so weiter. Jede hat eine eigene Einheit, jede hat einen eigenen Wert,
jede scheint unabhängig zu sein. Spätestens beim Studierenden taucht die
Frage auf: \emph{Warum gerade diese Werte? Warum nicht andere?}

Die übliche Antwort lautet, etwas verkürzt: ``Das wissen wir nicht. So ist
es eben.''

Diese Antwort ist intellektuell unbefriedigend. Sie suggeriert, die
Natur sei mit einem Dutzend unabhängig wählbarer Stellschrauben ausgestattet,
und der Mensch habe sie nachzumessen, ohne sie verstehen zu können. Schon
Sommerfeld, Eddington, Dirac und Feynman haben das nicht akzeptiert; jeder
auf seine Weise hat versucht, hinter die Mauer der ``fundamentalen Konstanten''
zu schauen.

FFGFT macht denselben Versuch, allerdings auf einem anderen Weg. Der Weg
beginnt mit einer scheinbar harmlosen Beobachtung: viele dieser ``Konstanten''
sind in Wirklichkeit keine. Sie sind \emph{Wechselkurse zwischen
Maßstabswahlen}, nicht eigenständige Naturgrößen. Das wird erst sichtbar,
wenn man konsequent in zwei Schritten reduziert.

\section{Erster Schritt: $\hbar = c = G = k_B = 1$}

Der erste Schritt ist Standard und in der theoretischen Physik seit
Jahrzehnten gängig: man arbeitet in \emph{natürlichen Einheiten}, in
denen die Konstanten $\hbar$, $c$, $G$ und $k_B$ einfach gleich Eins
gesetzt werden (Dok.~014, 015).

Was zunächst nach Notationsbequemlichkeit aussieht, ist tatsächlich eine
strukturelle Aussage. In SI-Einheiten haben Länge (Meter), Zeit (Sekunde)
und Masse (Kilogramm) unabhängige Dimensionen. Diese Unabhängigkeit ist
aber eine \emph{Konvention}, nicht eine Eigenschaft der Natur. Setzt man
$c = 1$, so werden Länge und Zeit zu derselben Größe. Setzt man $\hbar = 1$,
so werden Energie und inverse Zeit zu derselben Größe. Setzt man auch noch
$G = 1$ und $k_B = 1$, so kollabieren \emph{alle} mechanischen,
thermodynamischen und gravitativen Größen auf eine einzige Dimension:
Energie.

\begin{center}
\adjustbox{max width=\textwidth}{%
\begin{tabularx}{\linewidth}{Xl}
\toprule
Größe & Dimension \\
\midrule
Energie, Masse, Temperatur, Frequenz, Impuls & $[E]$ \\
Länge, Zeit, Wellenlänge & $[E]^{-1}$ \\
Kraft & $[E]^{2}$ \\
Geschwindigkeit & dimensionslos \\
Elektrische Ladung & dimensionslos \\
\bottomrule
\end{tabularx}%
}
\end{center}

Diese Tabelle (ausführlicher in Dok.~015) sagt eine einfache, aber wichtige
Sache: \emph{Die scheinbare Vielfalt der physikalischen Größen ist zum
großen Teil ein Artefakt der Maßstabswahl}. Es gibt nur eine fundamentale
Dimension, und alle übrigen sind Potenzen davon. Die ``vielen Konstanten''
der Schulphysik sind dann nicht mehr unabhängig --- sie sind
Umrechnungsfaktoren zwischen verschiedenen Energie-Potenzen.

Dieser Schritt allein ist nicht spezifisch für FFGFT. Er gehört zur
Standard-Toolkit der Hochenergiephysik. Aber er ist die Voraussetzung
für den zweiten, weniger gewohnten Schritt.

\section{Zweiter Schritt: $\alpha = 1$}

Auch in natürlichen Einheiten mit $\hbar = c = G = k_B = 1$ bleibt eine
Größe scheinbar fundamental, scheinbar unerklärt, scheinbar
``magisch'': die Feinstrukturkonstante
\[
  \alpha \;=\; \frac{e^{2}}{4\pi\varepsilon_{0}\hbar c} \;\approx\; \frac{1}{137{,}036}.
\]

Feynman hat sie als ``das größte Mysterium der Physik'' bezeichnet: eine
dimensionslose Zahl, die scheinbar aus dem Nichts kommt und doch die gesamte
Chemie und Atomphysik bestimmt.

FFGFT macht hier einen entscheidenden zweiten Schritt: \emph{auch $\alpha$
darf gleich Eins geschrieben werden}. Das ist keine reine Einheitenwahl
--- $\alpha$ ist dimensionslos und ändert sich bei einer Einheitenwahl
nicht, weil sich $e$ und $\varepsilon_{0}$ gemeinsam transformieren ---,
sondern eine Umdeutung, die an die in der Physik als
\textbf{Heaviside-Lorentz-Konvention} bekannte Schreibweise anschließt und
in Dok.~044 ausführlich diskutiert wird:

\begin{itemize}\setlength\itemsep{0pt}
  \item Im SI-System ist $\alpha = e^{2}/(4\pi\varepsilon_{0}\hbar c) \approx 1/137$.
  \item Im Gauss-System ist $\alpha = e^{2}/(\hbar c) \approx 1/137$.
  \item Im Heaviside-Lorentz-System (mit $\varepsilon_{0} = 1$
        und passender Wahl der Ladungsdefinition) wird $\alpha = e^{2}/(4\pi)$,
        mit demselben Zahlenwert $\approx 1/137$.
  \item Mit $e := \sqrt{4\pi\varepsilon_{0}\hbar c}$ wird die Kopplung
        in die Normierung des Feldes genommen: im Wechselwirkungsterm
        erscheint sie als Eins ($\alpha = 1$), vor dem Feldterm steht
        $1/(4\pi\alpha)$.
\end{itemize}

Alle diese Schreibweisen beschreiben dieselbe Physik. Was sich
ändert, ist nur, wo die Information sitzt --- ob in $\alpha$, in der
Ladung $e$, in der elektrischen Feldkonstante $\varepsilon_{0}$ oder in
der Normierung des Feldes.

Erst wenn man das durchschaut hat, wird die eigentliche Pointe sichtbar:
\emph{Es bleibt am Ende genau eine einzige dimensionslose Größe übrig, die
sich nicht durch Maßstabswahl wegeichen lässt.} In der FFGFT ist diese
Größe der geometrische Parameter
\[
  \xi_{0} \;=\; \frac{4}{30000} \;\approx\; 1{,}33 \times 10^{-4}.
\]

Alles andere --- $c$, $\hbar$, $G$, $k_B$, $\alpha$, und auch die Massen
der Teilchen, soweit sie als reine Verhältnisse zueinander betrachtet werden ---
folgt aus $\xi_{0}$ (bei den Massenverhältnissen zusammen mit den
gesetzten rationalen Vorfaktoren der Leiter), oder ist eine
Skalenkonvention.

\subsection*{Wichtige Klarstellung: ``auf 1 gesetzt'' heißt ``ausgeklammert''}

Die ``auf 1 gesetzten'' Werte $c, \hbar, G, k_B, \alpha$ verschwinden
nicht aus der Physik. Sie sind \emph{temporär ausgeklammerte
Skalenfaktoren}, keine wegrationalisierten Größen.

Solange man in \emph{Verhältnissen} rechnet --- Massenverhältnisse,
Frequenzverhältnisse, dimensionslose Kopplungen ---, heben sich diese
Faktoren heraus, und keine numerische Einsetzung ist nötig. Die Antwort
ist scharf, ohne dass $c$ oder $\hbar$ überhaupt erscheinen müssen.

Sobald man aber einen \emph{projizierten Messwert} in SI-Einheiten
erzeugen will --- die Elektronenmasse in Kilogramm, eine Frequenz in
Hertz, eine Energie in Joule ---, müssen die auf 1 gesetzten Werte
zusammen mit der fraktalen Korrektur \emph{wieder eingesetzt} werden.
Sie sind also bedingte Brückenfaktoren: abwesend in der reinen
Geometrie, präsent in der Projektion. Das ist kein Widerspruch, sondern
die saubere Form der Architektur: wer in der Geometrie bleibt, braucht
sie nicht; wer einen Messwert sehen will, holt sie zurück.

\section{Die Einsicht: zwei Schichten}

Aus diesen beiden Schritten ergibt sich die methodische Architektur, die
diesem Dokument seinen Namen gibt. Die FFGFT operiert auf \emph{zwei klar
getrennten Schichten}:

\begin{description}
  \item[Schicht~1 --- die Verhältnisse.] Alles, was als reines Verhältnis
    zwischen zwei gleichartigen Größen ausgedrückt werden kann --- also
    Massenverhältnisse wie $m_{\mu}/m_{e}$, Frequenzverhältnisse,
    Wicklungszahlen auf dem Torus, dimensionslose Kopplungen, der Parameter
    $\xi_{0}$ selbst --- ist \emph{einheiteninvariant}. Diese Verhältnisse
    sind dasjenige, was unabhängig von jeder Maßstabswahl ``wirklich da''
    ist. Die FFGFT macht ihre fundamentalen Aussagen auf dieser Schicht.

  \item[Schicht~2 --- die Umrechnung in SI.] Sobald man fragt, wie viele
    Kilogramm eine bestimmte Masse hat, wie viele Sekunden eine bestimmte
    Zeit dauert, oder wie viele Volt eine bestimmte Spannung trägt, betritt
    man eine zweite Schicht: die Übersetzung zwischen den dimensionslosen
    Verhältnissen der Schicht~1 und den menschlich gewählten Einheiten
    Kilogramm, Sekunde, Meter, Volt. Diese Übersetzung erfordert eine
    \emph{Skalenbrücke}, und in dieser Brücke --- und nur in dieser Brücke
    --- sitzen die fraktalen Korrekturen.
\end{description}

Das ist die zentrale, oft übersehene Aussage:

\begin{center}
\fbox{\parbox{0.85\textwidth}{\centering
  \textbf{Die fraktale Geometrie steckt in den Verhältnissen selbst ---
  aber dort als abgeschlossene, eingefrorene Rekursion.\\[4pt]
  Sichtbar aktiv wird sie erst in der Brücke zur SI-Sprache,
  wo die Einbettungskorrektur explizit werden muss.}
}}
\end{center}

Die Verhältnisse der Schicht~1 sind nicht fraktal-frei --- sie
\emph{sind} Fixpunkte fraktaler Rekursionen, die zur Ruhe gekommen
sind. Was fehlt, ist nicht die Fraktalität, sondern ihre sichtbare
Bewegung. Erst wer fragt, \emph{wie viel Kilogramm} eine Masse hat,
bricht den Fixpunkt auf und muss die Einbettungskorrektur explizit
mitführen. Dieser Übergang --- nicht die Physik selbst, aber die
Repräsentation der Physik in menschlichen Einheiten --- zahlt den
``Preis'', von dem Dok.~185 (\emph{Die Zeit als Einbettungspreis der
Mathematik}) spricht.

\section{Wo die Rekursion versteckt ist}

An dieser Stelle kann ein Missverständnis entstehen. Wenn Schicht~1
die reinen Verhältnisse enthält, scharf und ohne Korrektur --- wo
bleibt dann die Rekursion, von der FFGFT als \emph{fraktal-geometrische}
Theorie spricht? Hat man sie in den Verhältnissen verloren?

Die Antwort ist: nein. Sie ist in den Verhältnissen aufbewahrt. Aber
sie sieht dort still aus.

Eine reine Zahl wie das Verhältnis $m_{\mu}/m_{e} \approx 206{,}768$
wirkt statisch. Sie \emph{ist} aber das Endprodukt einer Rekursion, die
in dieser Zahl zur Ruhe gekommen ist. So wie der goldene Schnitt
$\varphi = (1+\sqrt{5})/2$ scheinbar eine schlichte Zahl ist, aber
tatsächlich die Grenze der Fibonacci-Folge $a_{n+1} = a_{n} + a_{n-1}$
--- die Rekursion sitzt nicht neben $\varphi$, sondern \emph{ist}
$\varphi$. So wie $\pi$ in den $n$-Ecken steckt und $e$ in der
Verzinsungsformel: eine reine Zahl ist immer ein eingefrorener Prozess.
Ein \emph{Fixpunkt} ist eine Geometrie, die ihre Bewegung in sich
aufgenommen hat.

Die FFGFT-Verhältnisse werden ebenso gelesen. Das Lepton-Verhältnis
$m_{\mu}/m_{e}$ gilt als Fixpunkt einer Wicklungsrekursion auf dem
4D-Torus. Die Exponenten der Leiter bilden dabei die geometrische Folge
$(3/2,\,1,\,2/3)$ mit dem Verhältnis $2/3$, keine Folge mit konstanter
Schrittweite (Dok.~210, W6). Die rationalen Vorfaktoren $4/3, 16/5, 25/9$
(Dok.~205) sind Setzungen der Leiter; ihre Herleitung aus einer
$(n,l,j)$-Formel trägt nicht. Als Deutung für die endliche Zahl der
Stufen führen Dok.~011 und 159 eine Fibonacci-artige Konvergenz an: die
Rekursionsamplitude $\varphi^{-2n}/\sqrt{5}$ erreicht bei $n \approx 8$
die Größenordnung von $\xi_{0}$, und dort schließt sich die Modenfolge
zu einer endlichen Tafel.

Die Rekursion ist in den Verhältnissen also vollständig anwesend --- aber
sie ist stillgelegt, weil sie konvergiert hat. Was wir als ``Verhältnis''
sehen, ist die finale Form einer Bewegung, die nun zur Ruhe gekommen
ist. Wer ein Verhältnis sieht, sieht eine fertige Geschichte. Wer fragt,
\emph{warum gerade dieses} Verhältnis, fragt nach der Rekursion, die zu
diesem Fixpunkt geführt hat. In FFGFT wird auf beide Ebenen geantwortet:
die Verhältnisse sind die Ergebnisse, die Wicklungen sind die
zugrundeliegende Geschichte.

\section{Wo die Rekursion sichtbar wird}

Wenn die Verhältnisse in sich abgeschlossen sind --- wo wird die
Rekursion dann \emph{in Bewegung} sichtbar? Antwort: in der Brücke zur
SI-Sprache.

Eine SI-Einheit (Kilogramm, Sekunde, Meter) ist nicht ein Fixpunkt; sie
ist eine \emph{willkürlich gewählte Zwischenstation} auf der Skalenleiter
der Natur. Die Planck-Skala und die elektroschwache Skala liegen rund 17
Größenordnungen auseinander; das Kilogramm sitzt irgendwo mittendrin,
festgelegt durch eine historisch gewachsene Maßstabswahl, nicht durch
einen geometrischen Eigenwert. Wer von einem geometrischen Eigenwert
in Schicht~1 zu einem Kilogramm-Wert in Schicht~2 übersetzt, fällt
nicht aus einer fertigen Rekursion in eine andere fertige, sondern
bewegt sich auf einer Skalentreppe, die in diesem Moment noch
\emph{in Bewegung} ist --- noch nicht konvergiert.

Genau dort taucht $K_{\mathrm{frak}} = 1 - 100\xi$ auf. Der Faktor 100
ist in Dok.~133 gesetzt, nicht hergeleitet; die dort gegebene Begründung
über den logarithmischen Abstand $\ln(10^{17}) \approx 39$ zwischen
Planck- und elektroschwacher Skala trägt nicht. Der Wert $1 - 100\xi$
wird durch den Vergleich mit $\alpha$ bestätigt, die Form von
$K_{\mathrm{frak}}$ ist offen. In der Lesart dieses Dokuments ist
$K_{\mathrm{frak}}$ die Spur einer Rekursion, die \emph{noch nicht zur
Ruhe gekommen ist}, weil unsere Maßstabswahl sie mittendrin abschneidet.

Die fraktale Geometrie zeigt sich also auf zwei Weisen:

\begin{itemize}\setlength\itemsep{2pt}
  \item in Schicht~1 als \emph{stillgelegte Rekursion} ---
        in den Verhältnissen, die Fixpunkte einer konvergierten
        Wicklungsfolge sind;
  \item in der Brücke als \emph{laufende Rekursion} ---
        in $K_{\mathrm{frak}}$, der kumulativen Spur einer
        Skalenkaskade, die wir mittendrin durch unsere Maßeinheiten
        unterbrechen.
\end{itemize}

\noindent
Beides ist dieselbe Geometrie. Einmal als Fixpunkt, einmal als
Skalenfluss. Was wie ``zwei verschiedene Phänomene'' aussieht, ist die
gleiche fraktale Selbstähnlichkeit in zwei Zuständen: am Ziel und
unterwegs.

\section{Konsequenzen}

Aus dieser Architektur folgen vier Konsequenzen, die das Verständnis vieler
einzelner FFGFT-Resultate spürbar vereinfachen.

\subsection{Verhältnis-Vorhersagen sind scharf}

Wenn die FFGFT eine Aussage über ein \emph{Verhältnis} macht --- zum
Beispiel über $m_{\mu}/m_{e}$ oder über das Verhältnis von Myon- zu
Tauon-Lebensdauer (Dok.~160) ---, dann ist diese Vorhersage scharf:
\emph{ohne} fraktale Korrektur, weil die Korrektur sich in einem Verhältnis
gleichartiger Größen herauskürzt.

\begin{center}
\fbox{\parbox{0.85\textwidth}{
  Wenn beide Größen im Zähler und Nenner durch denselben Umrechnungsfaktor
  zur SI-Skala übergehen, hebt sich dieser Faktor heraus. Was bleibt, ist
  reine Geometrie.
}}
\end{center}

Daher sind Verhältnis-Vorhersagen der FFGFT typisch auf zwei bis drei
Stellen genau: $m_\mu/m_e = (16/5)/(4/3)\cdot\xi^{-1/2} = 207{,}85$ gegen
gemessen $206{,}77$ ($+0{,}52\,\%$), $m_\tau/m_\mu =
(25/9)/(16/5)\cdot\xi^{-1/3} = 16{,}99$ gegen $16{,}82$ ($+1{,}0\,\%$).
Sie testen \emph{die Geometrie}, nicht die Umrechnung.

\subsection{SI-Werte tragen die fraktale Signatur}

Sobald die FFGFT eine Aussage über einen \emph{absoluten SI-Wert} macht ---
zum Beispiel über die Masse des Elektrons in Kilogramm, oder über die
Gravitationskonstante in m\textsuperscript{3}/(kg$\cdot$s\textsuperscript{2})
---, taucht die fraktale Korrektur auf. Das ist kein Defizit der Theorie,
sondern eine notwendige Konsequenz der Architektur:

\begin{itemize}\setlength\itemsep{0pt}
  \item Dok.~014 führt die Elektronenmasse als Zahl
        $m_{e}^{\,\mathrm{T0}} = 0{,}511$; das ist ihr Wert in MeV.
  \item Der Skalierungsfaktor $S_{T0} \approx 1{,}782662\times 10^{-30}$~kg
        ist genau 1~MeV/$c^2$ in kg
        ($10^6\cdot1{,}602177\times10^{-19}\,\text{J}/c^2 = 1{,}7826619\times10^{-30}$~kg);
        er ist nicht aus der fraktalen Geometrie abgeleitet.
  \item Das Produkt $0{,}511\cdot S_{T0} = 9{,}109\times10^{-31}$~kg
        ($m_{e}^{\,\mathrm{SI}} \approx 9{,}1094 \times 10^{-31}$~kg) ist
        damit die Einheitenumrechnung MeV$\to$kg.
\end{itemize}

Die fraktale Korrektur $K_{\mathrm{frak}} = 1 - 100\xi \approx 0{,}9867$
(Dok.~133) gehört in diese Brücke, nicht in die Elektronenmasse selbst.

\subsection{$D_f$ ist eine Eigenschaft der Brücke}

Daraus folgt eine wichtige Klärung zur fraktalen Dimension
$D_{f} \approx 2{,}999867$: sie sitzt nicht ``in der Masse'', auch nicht
``im Raum'' im naiven Sinne, sondern in der Brücke zwischen der
geometrisch-natürlichen Sprache und der SI-Sprache.

\begin{itemize}\setlength\itemsep{0pt}
  \item Die \emph{lokale} fraktale Dimension der Raumzeit
        $D_{f}^{\,\text{Raum}} = 3 - \xi \approx 2{,}9999$ ist eine
        Eigenschaft des Vakuums und liefert nur Korrekturen der
        Größenordnung $10^{-5}$, also praktisch Eins.
  \item Die \emph{effektive} fraktale Dimension
        $D_{f}^{\,\text{eff}} \approx 2{,}973$, die in der Brücke auftritt,
        ist eine \emph{kumulative} Größe: sie entsteht aus der Wirkung
        der lokalen Abweichung $\xi$ über den vollen
        Renormierungsgruppen-Lauf von der Planck-Skala zur elektroschwachen
        Skala (Dok.~133).
\end{itemize}

Beide haben denselben geometrischen Ursprung; aber sie sitzen an
verschiedenen Stellen der Architektur, und sie dürfen nicht verwechselt
werden.

\subsection{Status von $K_{\mathrm{frak}}$}

Eine häufige Frage an die FFGFT lautet: ``Ist $K_{\mathrm{frak}} = 1 -
100\xi$ nicht einfach so gewählt, dass die experimentellen Werte
herauskommen?'' Die Schichten-Architektur ordnet die Antwort:

\begin{itemize}\setlength\itemsep{0pt}
  \item $K_{\mathrm{frak}}$ gehört in die Brücke zwischen Schicht~1
        und Schicht~2, nicht in die Verhältnisse.
  \item Der Faktor $100$ ist in Dok.~133 gesetzt, nicht hergeleitet;
        die Begründung über den logarithmischen Abstand zwischen Planck-
        und elektroschwacher Skala ($\ln(10^{17}) \approx 39$) trägt
        nicht. Die Form von $K_{\mathrm{frak}}$ ist offen.
  \item Der Wert $1 - 100\xi$ wird durch den Vergleich der beiden Wege
        zu $E_0$ (nacktes und fraktal korrigiertes Mittel aus $m_e$ und
        $m_\mu$) mit $\alpha$ bestätigt; ein kleiner Rest bleibt offen.
        Eine $\alpha$-freie Messung von $K_{\mathrm{frak}}$ ist dieser
        Vergleich nicht.
\end{itemize}

Freier Parameter der Geometrie bleibt $\xi_{0}$; in Schicht~2 tritt der
eine Anker für den SI-Anschluss hinzu.

\section{Das Bild}

In Worte gefasst sagt die Schichten-Architektur folgendes:

\begin{center}
\fbox{\parbox{0.85\textwidth}{\centering
  \textbf{Die Natur rechnet in Verhältnissen.}\\[2pt]
  \textbf{Wir messen in SI.}\\[2pt]
  \textbf{Die fraktale Geometrie sitzt im Übersetzer.}
}}
\end{center}

Dies ist nicht nur eine pädagogische Faustregel, sondern eine
strukturelle Aussage über die FFGFT. Sie erklärt, warum manche
Verhältnis-Vorhersagen ohne fraktale Korrektur auskommen und andere eine
$1{,}3\%$-Korrektur erhalten; sie erklärt, wo $K_{\mathrm{frak}}$ sitzt; sie
erklärt, warum die Theorie mit einem einzigen freien Parameter
auskommt; und sie macht die Beziehung zwischen geometrischer
Idealisierung und experimenteller Messung präzise.

Die epistemische Pointe lautet: \emph{Was wir ``absolute Werte in SI''
nennen, sind menschliche Artefakte}. Das Universum ``weiß'' nichts vom
Meter, von der Sekunde oder vom Kilogramm. Es weiß nur von Verhältnissen,
von Wicklungszahlen, von Phasen. Die fraktalen Korrekturen sind der Preis,
den wir dafür zahlen, dass wir das Universum \emph{durch unsere willkürlichen
Einheiten hindurch} messen. Sie sind kein Defizit der Welt, sondern eine
Eigenschaft der Sprache, in der wir die Welt beschreiben.

\section{Die tiefere Pointe: das Universum rechnet nicht}

Selbst die Formulierung ``Schicht~1 \emph{berechnet} Verhältnisse,
Schicht~2 \emph{berechnet} die SI-Werte'' ist noch zu mechanistisch.
Tatsächlich findet in keiner der beiden Schichten ein Rechenvorgang
statt. Das Universum berechnet nichts; es \emph{ist} die Konfiguration.

Schicht~1 ist eine \emph{Geometrie zulässiger Modenkonfigurationen}.
Eine Mode existiert oder existiert nicht. Sie wird nicht ``errechnet'',
sie ist die geometrische Konfiguration selbst. Die Verhältnisse stehen,
wo sie stehen, weil die Wicklungsrekursion an dieser Stelle konvergiert
ist --- nicht weil irgendjemand sie ausgerechnet hätte.

Schicht~2, die Brücke, ist ebenfalls keine Rechnung, sondern eine
\emph{Übersetzung}. $K_{\mathrm{frak}}$ ist nicht ausgerechnet, sondern
Teil dieser Brücke. Auch die Skalentreppe
``rechnet'' nicht; sie ist da.

Berechnet wird nur dort, wo ein Beobachter Schicht~2 aus Schicht~1
ableiten will --- die Berechnung findet im \emph{Beobachter} statt,
nicht in der Welt.

Dieselbe Logik macht photonische Rechenmedien (TFLN-Kristalle,
Ising-Maschinen, analoge Resonatoren, Quantenpunkte; vgl. Dok.~161,
167, 173) so effizient: das Medium ``rechnet'' nicht, es liegt in der
Lösung. Die Geometrie der zugelassenen Phasenkonfigurationen \emph{ist}
das Rechenergebnis. Was bei einer digitalen Maschine als sequentielle
Operation in Zeit abläuft, ist im photonischen Medium eine
\emph{gleichzeitige} Konfiguration im Raum: nicht Algorithmus, sondern
Geometrie. FFGFT macht dieselbe Aussage über die Natur als Ganzes ---
die Modenstruktur des 4D-Torus ist nicht das Programm, das die Natur
ausführt, sondern die Natur selbst.

\section{Wo wir wirklich sind: c als Erfahrung der Vergänglichkeit}

Damit ist die Architektur aber noch nicht vollständig. Schicht~1 ist
die geometrische Reinform; Schicht~2 ist die menschliche
Maßstabsbeschreibung. Beides sind \emph{Sprachen} über etwas, in dem
wir tatsächlich sind. Wir leben aber weder in der reinen Geometrie noch
in den SI-Werten. Wir leben in einer dritten Lage, die in keiner der
beiden Schichten aufgeht.

Was wir tatsächlich erfahren, ist \emph{Vergänglichkeit}. Für uns ist
$c$ nicht 1. $c$ hat einen endlichen Wert, weil wir den Lichtweg als
Dauer erleben, als Strecke, die Zeit kostet. Nicht weil $c$ ``wirklich''
$299\,792\,458$~m/s wäre --- in der reinen Geometrie ist $c = 1$ ---,
sondern weil wir innerhalb der Konfiguration sitzen und sie als
\emph{Verstreichen} erfahren. Wären wir in Schicht~1 selbst, wären wir
nicht mehr Beobachter mit Lebenszeit, sondern die Geometrie. Die
auf~1 gesetzten Werte gehören in das Auge der Geometrie. Wir sehen aus
dem Inneren --- und genau dort haben sie endliche Werte, und genau dort
werden sie eingesetzt, sobald wir messen.

Das ist der Sinn der Aussage, dass die auf~1 gesetzten Werte bei
projizierten Messwerten wieder einzusetzen sind: nicht weil die
Geometrie sie hätte, sondern weil \emph{wir, die messen}, in einer Lage
sind, in der sie endliche Werte annehmen. $c$ ist die phänomenologische
Spur unserer Innenposition. Ohne Beobachter, der altert, gäbe es kein
endliches $c$.

\section{Schluss-Pointe: die fraktale Spur unseres Erlebens}

Und damit schließt sich die Architektur, die zu Beginn nur statisch
aussah, zu einer dynamischen Form. Die Rekursion war nie nur in den
Verhältnissen (Schicht~1) oder nur in der Brücke (Schicht~2). Sie ist
auch in uns.

Jeder Moment, in dem wir altern, enthält alle vorherigen Momente in
gedämpfter Form. Unser Erleben ist nicht ein linearer Strom, sondern
ein \emph{Konzertsaal} (Dok.~159): jeder neue Klang trägt das Echo
aller alten in sich. Das menschliche Gedächtnis ist nicht ein Speicher,
sondern eine rekursive Resonanz, in der jede Wiederkehr eine leise
Kopie der vorherigen ist.

Wenn das Universum die Konfiguration \emph{ist} und wir die Konfiguration
\emph{erleben}, dann ist unser Erleben die \emph{fraktale Spur der
Geometrie, in der wir liegen}. Die Rekursion ist nicht ein abstraktes
Merkmal der Theorie --- sie ist dasjenige, woran wir uns als Beobachter
wiedererkennen. Das gleiche Muster, das in den Verhältnissen zur Ruhe
gekommen ist und in der Skalenbrücke noch in Bewegung ist, ist auch
das Muster unseres Wahrnehmens.

\medskip

Das ist der Bogen, den dieses Dokument schlägt: von der scheinbaren
Vielheit der Naturkonstanten über die zwei Schichten und ihre rekursive
Beziehung bis zur Position des Beobachters, der die Geometrie nicht
\emph{kennt}, sondern \emph{ist} --- aber sie als Vergänglichkeit
durchläuft. Die FFGFT beschreibt keine Maschinerie, die jemand von
außen anschaut. Sie beschreibt eine Geometrie, die sich selbst trägt,
und einen Beobachter, der innerhalb dieser Geometrie liegt und sie
deshalb in zwei Sprachen ausdrücken muss.

Wer das gesehen hat, sieht $c$, $\hbar$, $\alpha$ und $K_{\mathrm{frak}}$
nicht mehr als getrennte Naturkonstanten, sondern als verschiedene
Spuren derselben einen Geometrie --- gesehen aus verschiedenen Lagen
der einen Konfiguration, die wir nicht beobachten, sondern in der wir
sind.

\section{Lese-Reihenfolge: wo welche Schicht ausgeführt ist}

Wer die FFGFT nun systematischer durcharbeiten möchte, findet die
beiden Schichten in folgenden Dokumenten ausführlich behandelt:

\begin{itemize}\setlength\itemsep{0pt}
  \item \textbf{Schicht~1 (Verhältnisse, geometrische Reinform):}
    Dok.~002 (Reise), Dok.~015 (Systematik natürlicher Einheiten),
    Dok.~009 ($\xi$-Ursprung), Dok.~205 (Narrativ),
    Dok.~006 (Teilchenmassen als Verhältnisse),
    Dok.~060/116/159 (Moden- und Wicklungsverhältnisse),
    Dok.~230 (Hilbertraum-Brücke, ebenfalls verhältnis-basiert).
  \item \textbf{Schicht~2 (Skalenbrücke, fraktale Korrekturen):}
    Dok.~014 (Übergang nat./SI mit Skalierungsfaktor $S_{T0}$),
    Dok.~013 (SI-Konventionen), Dok.~133 (Vollständige Herleitung von
    $K_{\mathrm{frak}}$), Dok.~134 (c als Konvention), Dok.~044
    (Konventionen für $\alpha$),
    Dok.~185 (Einbettungspreis der Mathematik).
\end{itemize}

Wer die Dokumente in dieser Reihenfolge liest, sieht die Architektur, von
der dieses Dokument spricht, an konkreten Stellen wirken: Verhältnisse
oben, Brücke unten, kein freier Parameter außer $\xi_{0}$.

\section*{Zusammenfassung}

\begin{itemize}\setlength\itemsep{2pt}
  \item FFGFT operiert auf zwei klar getrennten Schichten:
        Verhältnisse (Schicht~1, fraktal-frei) und SI-Projektion
        (Schicht~2, mit fraktaler Brücke).
  \item Die ``auf~1 gesetzten'' Werte $c, \hbar, G, k_B, \alpha$ sind
        nicht weg, sondern ausgeklammert; sie kehren zurück, wenn man
        projizierte Messwerte erzeugt.
  \item Die Rekursion ist in den Verhältnissen als stillgelegte
        Konvergenz anwesend (Fixpunkte einer Wicklungsfolge mit
        Fibonacci-Konvergenz bei $n \approx 8$) und in der Brücke als
        laufende Skalenkaskade ($K_{\mathrm{frak}} = 1 - 100\xi$).
  \item Das Universum berechnet nicht; es ist die Konfiguration.
        Berechnung findet im Beobachter statt --- analog zur Logik
        photonischer Rechenmedien (TFLN, Ising-Maschinen).
  \item Wir leben weder in Schicht~1 noch in Schicht~2, sondern in
        einer dritten Lage: der Erfahrung der Vergänglichkeit. $c$ hat
        für uns einen endlichen Wert, weil wir die Konfiguration
        \emph{durchlaufen}, statt sie zu \emph{sein}.
  \item Die Rekursion zeigt sich auch in unserem Erleben: jeder Moment
        trägt das Echo aller vorherigen. Unser Wahrnehmen ist die
        fraktale Spur der Geometrie, in der wir liegen.
\end{itemize}

\medskip
\noindent\textit{Stand: Mai 2026. Dieses Dokument ist
methodisch-pädagogisch und enthält keine eigenständigen physikalischen
Aussagen; alle Sachaussagen verweisen auf die jeweils genannten
Quelldokumente.}
